\section{Descripción de Criterios de Selección}\label{Ap:Crit}

Mediante el Proceso de Análisis Jerárquico (AHP), el cual es un método de selección multi-criterio se procedió a evaluar las distintas propuestas de solución que buscan solucionar todas las necesidades identificadas previamente.

Para poder realizar la selección de los diferentes componentes se consideraron una serie de criterios, los cuales son los siguientes:

\begin{itemize}
    \item Costo:
    Mientras el precio sea menor se le asigna una mayor calificación.
    \item Facilidad de manufactura:
    Se refiere a que tan fácil es manufacturar, en caso de ser necesario, el componente a utiliza.
    \item Disponibilidad:
    Es la facilidad con la que es posible conseguir algún componente en el mercado.
    \item Seguridad de uso:
    Representa el riesgo que puede presentar para una persona determinado sistema o componente, mientras más peligroso menor ponderación obtiene.
    \item Facilidad de uso:
    Mientras más sencillo sea de utilizar por personas sin los conocimientos técnicos, mejor evaluación.
    \item Integración:
    Hace referencia a la facilidad de interconectar un componente con otro, se busca tener un alto nivel de integración de un componente con otro.
    \item Durabilidad:
    Se refiera la tiempo de vida útil de cada componente.
    \item Mantenimiento:
    Facilidad con la cual pueda realizarse el mantenimiento de cierto componente, mientras más fácil mejor calificación.
    \item Ergonomía:
    Se trata de la relación y comodidad que se genera al utilizar algún componente, se busca una alta ergonomía.
    \item Portabilidad:
    Facilidad de desplazar un componente de un lugar a otro, se busca una alta portabilidad.
    \item Estética
    Agrado visual que genera cierto componente, sin tener en cuenta su funcionabilidad.
\end{itemize}

Estos diferentes criterios fueron ponderados, utilizando el criterio AHP, en una tabla, como se muestra en el apéndice \ref{Ap:Crit}, para determinar qué características son las más sobresalientes y que tendrán mayor peso al momento de realizar la selección de los elementos del sistema mecatrónico.

Aquí se hace la presentación de los valores utilizados para la evaluación de las diferentes soluciones tecnológicas propuestas, para elegir las mejores opciones.

\begin{table}
    \centering
    \caption{Cuantificación de los criterios}
    \resizebox{!}{20cm}{
    \begin{turn}{90}
      \begin{tabular}{|m{1.5cm}|c|m{1.5cm}|m{1.5cm}|m{1.5cm}|m{1.5cm}|m{1.5cm}|m{1.5cm}|m{1.5cm}|m{1.5cm}|m{1.5cm}|m{1.5cm}|m{1.5cm}|c|c|} \hline 
         Criterio&  &  Costo&  Facilidad de manufactura&  Disponi- bilidad&  Seguridad de uso&  Facilidad de uso&  Integra- cion&  Durabili- dad&    Manteni- miento&Ergono- mia&Portabi- lidad&Estetica  &Suma&Porcentaje\\ \hline 
         Costo&  1& \cellcolor[rgb]{ 1,  1,  0} X&  1&  1&  2&  2&  2&  2&    2&2&2&  2&18&16.7 \%\\ \hline 
         Facilidad de manufactura&  2&  1& \cellcolor[rgb]{ 1,  1,  0} X&  2&  2&  1&  1&  2&    2&2&2&  2&17&15.7 \%\\ \hline 
         Disponi- bilidad&  3&  1&  0& \cellcolor[rgb]{ 1,  1,  0} X&  1&  1&  2&  2&    2&2&2&  2&15&13.9 \%\\ \hline 
         Seguridad de uso&  4&  0&  0&  1& \cellcolor[rgb]{ 1,  1,  0} X&  2&  2&  2&    2&2&2&  2&15&13.9 \%\\ \hline 
         Facilidad de uso&  5&  0&  1&  1&  0& \cellcolor[rgb]{ 1,  1,  0} X&  2&  1&    1&1&2&  2&11&10.2 \%\\ \hline 
         Integra- ción&  6&  0&  1&  0&  0&  0& \cellcolor[rgb]{ 1,  1,  0} X&  0&    2&2&2&  2&9&8.3 \%\\ \hline 
         Durabili- dad&  7&  0&  0&  0&  0&  1&  0& \cellcolor[rgb]{ 1,  1,  0} X&    1&1&2&  2&7&6.5 \%\\ \hline 
         Manteni- miento&  8&  0&  0&  0&  0&  1&  0&  1&  \cellcolor[rgb]{ 1,  1,  0}  X&2&2&  2&8&7.4 \%\\ \hline 
         Ergono- mia&  9&  0&  0&  0&  0&  1&  0&  1&    0& \cellcolor[rgb]{ 1,  1,  0} X&2&  2&6&5.6 \%\\ \hline 
 Portabi- lidad& 10& 0& 0& 0& 0& 0& 0& 0&   0&0& \cellcolor[rgb]{ 1,  1,  0} X&  1&1&0.9 \%\\ \hline 
 Estetica& 11& 0& 0& 0& 0& 0& 0& 0&   0&0&0& \cellcolor[rgb]{ 1,  1,  0} X &1&0.9 \%\\ \hline 
 Total& & & & & & & & &   &&&  &108&100.00\%\\ \hline
    \end{tabular}  
    \end{turn}
    }
    \label{tab:tabcri}
\end{table}